% =====================================================================
%  Rubrica de evaluacion de calidad de requisitos funcionales
%  Proyecto SIMPA -- Equipo AHMRV -- ISR-401 -- UTEQ -- 2026-2027 PPA
%  Instrumento del cuasi-experimento del Enfoque 1
%  Compilacion: pdflatex rubrica_evaluacion.tex  (x2)
% =====================================================================
\documentclass[11pt,a4paper]{article}

\usepackage[utf8]{inputenc}
\usepackage[T1]{fontenc}
\usepackage[margin=2.3cm]{geometry}
\usepackage{array}
\usepackage{booktabs}
\usepackage{longtable}
\usepackage[table]{xcolor}
\usepackage{enumitem}
\usepackage{tcolorbox}
\usepackage{fancyhdr}
\usepackage{titlesec}
\usepackage{lastpage}
\usepackage[hidelinks]{hyperref}

\definecolor{uteqverde}{HTML}{1F6F3C}
\definecolor{grisclaro}{HTML}{F2F2F2}

\newcolumntype{L}[1]{>{\raggedright\arraybackslash}p{#1}}
\newcolumntype{C}[1]{>{\centering\arraybackslash}p{#1}}

\titleformat{\section}{\normalfont\large\bfseries\color{uteqverde}}{\thesection.}{0.5em}{}
\titlespacing*{\section}{0pt}{12pt}{5pt}
\titleformat{\subsection}{\normalfont\normalsize\bfseries}{\thesubsection}{0.5em}{}

\pagestyle{fancy}
\fancyhf{}
\renewcommand{\headrulewidth}{0.4pt}
\lhead{\footnotesize\color{uteqverde}Rúbrica de evaluación de requisitos}
\rhead{\footnotesize SIMPA --- Equipo AHMRV}
\cfoot{\footnotesize Página \thepage\ de \pageref{LastPage}}

\setlength{\parskip}{4pt}
\setlength{\parindent}{0pt}

\newcommand{\dimtitulo}[1]{\subsection*{\color{uteqverde}#1}}

\begin{document}

\begin{center}
{\footnotesize UNIVERSIDAD TÉCNICA ESTATAL DE QUEVEDO\\
Facultad de Ciencias de la Computación --- Carrera de Software}\\[10pt]
{\Large\bfseries\color{uteqverde} RÚBRICA DE EVALUACIÓN}\\[3pt]
{\large Calidad de requisitos funcionales}\\[2pt]
{\footnotesize Instrumento del estudio registrado en OSF \texttt{4z35d}}
\end{center}

\vspace{6pt}

\section*{Qué se le pide}

Usted va a leer \textbf{50 requisitos funcionales} de un sistema de gestión de
cultivo de palma africana y puntuar cada uno en \textbf{cinco dimensiones de
calidad}, con una escala de 1 a 5.

\begin{tcolorbox}[colback=grisclaro,colframe=uteqverde,boxrule=0.6pt,arc=2pt]
\textbf{Tres condiciones del estudio:}
\begin{itemize}[nosep,leftmargin=1.4em]
  \item Los requisitos proceden de \textbf{dos fuentes distintas}, y usted no
        sabe cuál corresponde a cada uno. Es deliberado: si lo supiera, su
        juicio podría verse afectado.
  \item Puntúe \textbf{de forma independiente}, sin comentar sus puntuaciones
        con las otras personas evaluadoras durante la sesión.
  \item \textbf{No se evalúa el sistema ni a sus autores}, sino la calidad de
        redacción de cada requisito como enunciado.
\end{itemize}
\end{tcolorbox}

\section{Escala común}

Las cinco dimensiones usan la misma escala:

\renewcommand{\arraystretch}{1.3}
\begin{longtable}{C{1.2cm} L{13.5cm}}
\toprule
\rowcolor{grisclaro}
\textbf{Nivel} & \textbf{Significado} \\
\midrule
\endhead
\textbf{1} & La propiedad está \textbf{ausente}. \\
\textbf{2} & Presente de forma \textbf{marginal}. \\
\textbf{3} & Presente con \textbf{deficiencias apreciables}. \\
\textbf{4} & Presente con \textbf{deficiencias menores}. \\
\textbf{5} & \textbf{Plenamente presente}. \\
\bottomrule
\end{longtable}

Si duda entre dos niveles, elija el \textbf{menor}. Es preferible una
puntuación conservadora y consistente que una generosa y variable.

\section{Las cinco dimensiones}

\dimtitulo{1. Completitud}

\textbf{Pregunta guía:} ¿el requisito contiene toda la información necesaria
para ser implementado, sin que haga falta consultar otra fuente?

Un requisito completo indica \emph{quién} realiza la acción, \emph{qué} hace el
sistema y \emph{bajo qué condiciones}. Si faltan datos que un desarrollador
tendría que adivinar, la dimensión está afectada.

\begin{tcolorbox}[colback=white,colframe=uteqverde!45,boxrule=0.5pt,arc=2pt]
\textbf{Nivel 5.} «El sistema debe permitir al supervisor registrar una labor
agrícola indicando lote, tipo de labor, fecha, cuadrilla asignada y avance en
hectáreas, y almacenarla asociada al lote seleccionado.»\\[4pt]
\textbf{Nivel 1.} «El sistema debe gestionar labores.»
\end{tcolorbox}

\dimtitulo{2. Ausencia de ambigüedad}

\textbf{Pregunta guía:} ¿el requisito admite una sola interpretación?

Términos como «rápido», «amigable», «adecuado», «si es necesario» o «entre
otros» introducen ambigüedad. También la introduce una conjunción que permita
leer dos comportamientos distintos.

\begin{tcolorbox}[colback=white,colframe=uteqverde!45,boxrule=0.5pt,arc=2pt]
\textbf{Nivel 5.} «El sistema debe mostrar la alerta en la pantalla de inicio
del usuario cuyo rol sea supervisor o administrador.»\\[4pt]
\textbf{Nivel 1.} «El sistema debe notificar oportunamente a los usuarios
correspondientes cuando sea adecuado.»
\end{tcolorbox}

\dimtitulo{3. Verificabilidad}

\textbf{Pregunta guía:} ¿se puede diseñar una prueba que determine, sin
discusión, si el requisito se cumple?

Un requisito verificable produce un resultado observable. Si para comprobarlo
hiciera falta un juicio subjetivo, la dimensión está afectada.

\begin{tcolorbox}[colback=white,colframe=uteqverde!45,boxrule=0.5pt,arc=2pt]
\textbf{Nivel 5.} «Al guardar un registro fitosanitario, el sistema debe
mostrarlo en la tabla de historial del lote seleccionado dentro de la misma
sesión.»\\[4pt]
\textbf{Nivel 1.} «El sistema debe ofrecer una buena experiencia de registro
fitosanitario.»
\end{tcolorbox}

\dimtitulo{4. Corrección respecto de la fuente}

\textbf{Pregunta guía:} ¿lo que el requisito afirma es coherente con el dominio
descrito en el material de contexto que se le ha entregado?

Se evalúa aquí si el requisito introduce elementos que el dominio no sostiene,
o si contradice lo descrito. Un requisito puede estar impecablemente redactado
y ser incorrecto.

\begin{tcolorbox}[colback=white,colframe=uteqverde!45,boxrule=0.5pt,arc=2pt]
\textbf{Nivel 5.} «El sistema debe registrar el conteo de flores polinizadas
por recorrido del trabajador.»\\[4pt]
\textbf{Nivel 1.} «El sistema debe calcular automáticamente el precio
internacional del aceite de palma para fijar el salario de la cuadrilla.»
\end{tcolorbox}

\dimtitulo{5. Consistencia interna}

\textbf{Pregunta guía:} ¿el requisito es coherente consigo mismo y con la
terminología que emplea?

Se evalúa si el enunciado usa el mismo término para el mismo concepto de
principio a fin, y si sus partes no se contradicen.

\begin{tcolorbox}[colback=white,colframe=uteqverde!45,boxrule=0.5pt,arc=2pt]
\textbf{Nivel 5.} «El sistema debe permitir al operario consultar las labores
asignadas a su cuadrilla, y solo esas labores.»\\[4pt]
\textbf{Nivel 1.} «El sistema debe permitir al operario consultar todas las
parcelas, restringiendo el acceso del operario a las parcelas de su lote.»
\end{tcolorbox}

\section{Antes de empezar}

\begin{enumerate}[leftmargin=1.6em]
  \item Lea esta rúbrica completa y plantee sus dudas \textbf{antes} de puntuar
        el primer requisito. Una vez iniciada la sesión, no se resuelven dudas
        de criterio, porque afectarían a la comparabilidad entre evaluadores.
  \item Puntúe los requisitos \textbf{en el orden en que aparecen} en el
        cuadernillo.
  \item Puntúe \textbf{las cinco dimensiones} de cada requisito antes de pasar
        al siguiente.
  \item Si un requisito le resulta imposible de puntuar en alguna dimensión,
        anótelo en la columna de observaciones en lugar de dejarlo en blanco.
\end{enumerate}

\begin{tcolorbox}[colback=grisclaro,colframe=uteqverde,boxrule=0.6pt,arc=2pt]
\textbf{Sobre la duración.} Son 50 requisitos por 5 dimensiones: 250
puntuaciones. Se estima entre 60 y 90 minutos. Puede tomar los descansos que
necesite. No hay ventaja en terminar antes.
\end{tcolorbox}

\vfill
{\footnotesize
Instrumento del protocolo registrado en OSF \texttt{4z35d}. Proyecto SIMPA ---
Equipo AHMRV --- Ingeniería de Requerimientos (ISR-401) --- Universidad Técnica
Estatal de Quevedo --- Período 2026--2027 PPA.}

\end{document}
